\section{Discussion}\label{sec:discussion}

Tables~\ref{tab:blv_cr}, ~\ref{tab:blv_tn}, and ~\ref{tab:blv_gg} show Cohen's $d$, which is the size of the effect of the treatment in the respective table. Ratings on \textbf{Nature} are not included in the average computation since it is a categorical variable; \textit{i.e.}, a low \textbf{Nature} rating simply means the description was perceived to be more straight facts-oriented than commentary-oriented, and not necessarily of a lower quality.

\paragraph{Combined Effect Size}
Table~\ref{tab:blv_cr} shows the effect size of fine tuning on \textsc{Sightation} \textit{and} applying the guided generation prompt at test time. With the combined recipe applied, the 2B model achieves an average of $0.36\sigma$ units of improvement, while the 7B model, $0.58\sigma$ units.
Intriguing observations can be made on succinctness. The 2B model exhibited the smallest effect size in this aspect, whereas the 7B model achieved the highest enhancement. This suggests that the combined recipe applied on the smaller model had negligible effect in making its descriptions more succinct. In fact, the combined recipe enhanced \textbf{Nature} by a large effect ($1.08\sigma$), implying that, with smaller models, the prime importance of the combined recipe lies in shaping the descriptions to be far more interpretive. The opposite can be said of the 7B model: the combined recipe greatly ($1.69\sigma$) enhances its succinctness, whilst shaping its descriptions far less interpretive ($-2.38\sigma$) and straight facts-oriented instead. This is in line with 3 separate comments by our BLV educators (B1, B2, and B5) who have, unknowingly of each other's interview responses, stressed the importance of succinctness: ``The description must deliver all visual items in an accurate and consistent manner, with not too long a text and including the key elements.''

\paragraph{Tuning Effect Size}
Table~\ref{tab:blv_tn} shows the effect size of fine tuning on \textsc{Sightation}. For instance, with guided generation absent, the 2B model still reaps $0.87\sigma$ units of improvement in the diversity aspect of its descriptions. The improvement margin is even amplified further by applying guided generation on the tuned model, except for usefulness in solving questions.
The table shares the observation made on the succinctness-nature relationship conveyed in Table~\ref{tab:blv_cr}, albeit to a lesser extent on the 7B model with guided generation. This set, whose ratings are on the rightmost column of Table~\ref{tab:blv_tn}, showed meaningful effect size only in usefulness as a summary and nature. This implies that larger models are already somewhat capable of capitalizing on the guided generation prompt at test time and carry less reliance on the fine tuning process.


\paragraph{Guided Generation Effect Size}
Table~\ref{tab:blv_gg} shows that the guided generation yields benefits even to GPT, possibly indicative of the under-representation of BLV accessibility needs and preferences in the pre-training data. It is important to note that, for the 2B model, the best effect of guided generation is achieved only \textit{after} the model is tuned on our dataset, again highlighting the BLV alignment capabilities of our dataset, that cannot be mimicked by test-time prompt engineering alone.
